\documentclass[12pt]{article}
\usepackage[margin=1in, paperwidth=8.5in, paperheight=11in]{geometry}
\usepackage{ifpdf,amsmath, amssymb, comment, color, graphicx, stmaryrd,setspace,enumerate,tikz, fancyhdr, wrapfig, textcomp, mathptmx, siunitx}
\usepackage{pdfpages}
\usepackage{enumerate}

\usepackage{hyperref}
\hypersetup{
    colorlinks=true,
    urlcolor=blue,
}

\setlength{\headheight}{14.5pt}
\newcommand{\Q}{\mathbb{Q}}
\newcommand{\R}{\mathbb{R}}
\newcommand{\Z}{\mathbb{Z}}

\newenvironment{red}{\color{red}}{\ignorespacesafterend}


\newcommand{\blue}[1]{\textcolor{blue}{#1}}
\newcommand{\green}[1]{\textcolor{green}{#1}}
\renewcommand{\section}[1]{\begin{center} \textbf{#1} \\\end{center}}
%
\setlength{\parindent}{0in}
%\oddsidemargin=-.25in
\allowdisplaybreaks
\pagestyle{fancy}
\renewcommand{\headrulewidth}{0pt}
\lhead{MATH 202}
\rhead{Fall 2024}
%\lfoot{\copyright\ CLEAR Calculus 2010}
\cfoot{}
\renewcommand{\thefootnote}{*} 
\hyphenpenalty=10000
\begin{document}
%


%\onehalfspacing
\allowdisplaybreaks
%##################################################################
\section{Physics applications of definite integrals}
Here is a quick reference for the answers to this problem set.

\hrulefill

\begin{enumerate}
\item The kinetic energy of the second hand is $\SI{5.483e-7}{J}$.

\item The total mass of the oil slick is $\SI{3.1386e6}{kg}.$

\item The energy to launch the CubeSat to the moon is $\SI{1.227e8}{J}$.

\item The energy to push the movable magnet is $k\cdot\SI{13.33}{J}$.

\item The total mass of the atmosphere is $\SI{5.162e18}{kg}$.

\item The gravitational attraction between the rod and the ball is $\dfrac{GMm}{a(a+\ell)} \si{N}$.
\end{enumerate}
%#####################################d############################
\end{document}